% ═════════════════════════════════════════════════════════════════════════════
\section{Deductive and Proof-Based Methods}
\label{sec:deductive}
% ═════════════════════════════════════════════════════════════════════════════

Deductive approaches shift the verification task from automated state-space
exploration to the construction of formal proofs, often assisted by theorem
provers or proof assistants. The idea is to find invariants (properties
that hold at every reachable state of the system for every~$n$) and to prove
that they are preserved by every system transition, regardless of the number
of processes~\cite{OG76,Lam94}. Since the proof is carried out symbolically
over an arbitrary~$n$, this style of reasoning handles unbounded parameters
without requiring enumeration of instances.

\subsection{Inductive Invariants}

We say that a set of configurations is an \emph{inductive invariant} if it
contains all initial configurations and is closed under the transition
relation, that is, every successor of a configuration in the set is again in
the set. An inductive invariant contains all reachable configurations, so if
it is disjoint from the bad configurations, the system is safe. In the
parameterized setting, invariants are expressed as formulas that quantify
over process indices, and the checks are carried out for an arbitrary
number~$n$ of processes.

For the running example, mutual exclusion is expressed by the formula stating
that for all distinct indices~$i$ and~$j$, processes~$i$ and~$j$ are not both
in~$\Crit$. This formula is itself an inductive invariant. It holds initially,
since all processes are in~$\Idle$. A step from~$\Idle$ to~$\Wait$ or
from~$\Crit$ to~$\Idle$ does not create a process in~$\Crit$. A step in which
some process~$k$ moves from~$\Wait$ to~$\Crit$ is only enabled when no other
process is in~$\Crit$, so after the step, $k$ is the only process in~$\Crit$.
The argument does not depend on~$n$ or on~$k$, and it is therefore a proof for
all instances at once.

In realistic protocols, the property to be proved is rarely inductive by
itself. If the universal guard of the running example is replaced by a loop in
which a process checks the other processes one at a time, then a process may
have checked some of the others while they were outside the critical section,
and these processes may have entered it since. The property must then be
\emph{strengthened} with auxiliary assertions, for instance relating the
progress of a process through its loop to the states of the processes it has
already checked. Finding such strengthenings is the creative part of a
deductive proof, and it is the step that the techniques in the rest of this
section, and in Section~\ref{sec:smt}, try to automate.

\subsection{Assertional Reasoning and TLA+}

Classic work by Owicki and Gries~\cite{OG76} laid the groundwork for
assertional reasoning about concurrent programs: each process is annotated
with pre- and post-conditions, and an interference freedom check ensures that
no process's execution invalidates another's assertions. The method was
designed for a fixed, known number of processes, but the underlying proof
discipline (inductive invariants together with noninterference) scales
naturally to the parameterized case by carrying out the annotation and check for a
symbolic, arbitrary process in a family.

Lamport's Temporal Logic of Actions~\cite{Lam94} and the specification
language TLA+~\cite{Lam02} provide a formal language for specifying and proving
properties of concurrent and distributed systems in this spirit. A system is
described by an initial predicate and a next-state relation, and parameters
such as the set of processes are left as constants of the specification. The
TLA+ proof system~\cite{CDLMRV12} supports machine-checked proofs of
invariants for all values of these constants, while the TLC model checker is
used to check the same specification for small, fixed values. Checking small
instances with TLC has proved effective in finding subtle errors in
industrial-scale protocols, for instance at Amazon Web
Services~\cite{NRZBMD15}, and machine-checked proofs extend such results to
all values of the parameters.

\subsection{Invisible Invariants}

A significant challenge for deductive methods is the discovery of inductive
invariants: properties strong enough to be preserved under all transitions and
to imply the desired safety property, but not so strong that they exclude
reachable configurations.
The \emph{invisible invariants} technique of Pnueli, Ruah, and
Zuck~\cite{PRZ01,APRXZ01} addresses this by automatically generating candidate
invariants from small finite instances of the system and then checking whether
they hold inductively for all~$n$. The idea is that the right invariant is
often already visible in small instances and can be generalized to the full
parameterized family.

The procedure works as follows. First, the set of reachable configurations of
an instance of a small size is computed symbolically, for example with BDDs.
Second, this set is projected onto one or two processes, which yields a
formula $\psi(i)$ or $\psi(i,j)$ describing the possible combinations of local
states of these processes. Third, the formula is generalized to the candidate
invariant ``for all~$i$, $\psi(i)$'' or ``for all distinct~$i$ and~$j$,
$\psi(i,j)$''. Finally, the candidate is checked to be inductive and to imply
the property. For a restricted class of systems, the last check can itself be
carried out on an instance of bounded size, by a \emph{small-model theorem}
that bounds the size of a counterexample to inductiveness. In that case the
whole procedure is automatic, and the invariant never needs to be seen by the
user, which explains the name~\cite{ZP04}.

In the running example, the reachable configurations of the instance of size
two, projected onto both processes, give the formula $\psi(i,j)$ stating that
$i$ and~$j$ are not both in~$\Crit$. Its generalization is the invariant from
the beginning of this section. This makes deductive verification partly
automatic, at the cost of completeness: if the right invariant is not among
those generated, the procedure does not terminate with a proof.

\subsection{Interactive Proofs}

More recent work integrates deductive reasoning with interactive tools,
allowing arbitrarily complex invariants to be stated and verified with
machine-checked proofs. The Ivy system~\cite{PMPSS16}, which we return to in
Section~\ref{sec:smt}, is representative of this direction. A
\emph{counterexample to induction} for a candidate invariant is a
configuration that satisfies the candidate but has a successor that does not.
When a candidate invariant is not inductive, Ivy presents such a
counterexample. The user generalizes this configuration into a pattern that should be
excluded, and the tool adds the corresponding conjunct to the invariant. The
checks are carried out automatically, so the interaction focuses on the
insight that the user brings, namely which features of a counterexample
explain why it is unreachable.

Proof-based methods tend to be the most expressive and general of all
approaches to parameterized verification, capable of handling protocols with
complex data structures, arithmetic, and dynamic topology. Their principal
limitation is the burden they place on the user: finding and expressing the
right invariant requires deep insight into the protocol's correctness
argument, and scaling to large or complex protocols remains demanding even
with modern tool support.
